\textbf{Phần 1. Trắc nghiệm lựa chọn}
\setcounterpageref{ex}{0}

\begin{ex} %[3D8N49]
Số gồm $5$ nghìn, $6$ trăm, $0$ chục và $4$ đơn vị là
\choice
{$5640$}
{$5064$}
{\True $5604$}
{$6504$}
\loigiai{
Số đó là $5604$.\\
\textbf{Đáp án của câu hỏi là C}
}
\end{ex}

\begin{ex} %[3D8N49]
Chữ số La Mã $\text{XVI}$ biểu thị số nào?
\choice
{$14$}
{\True $16$}
{$15$}
{$17$}
\loigiai{
$\text{XVI} = 10 + 6 = 16$.\\
\textbf{Đáp án của câu hỏi là B}
}
\end{ex}

\begin{ex} %[3D8H49]
Điền dấu thích hợp: $7\,450 \dots 7\,540$.
\choice
{$>$}
{$=$}
{Không so sánh được}
{\True $<$}
\loigiai{
Hàng trăm: $4 < 5$ nên $7450 < 7540$.\\
\textbf{Đáp án của câu hỏi là D}
}
\end{ex}

\begin{ex} %[3D8H49]
Làm tròn số $4\,785$ đến hàng trăm ta được
\choice
{$4700$}
{$4790$}
{\True $4800$}
{$5000$}
\loigiai{
Hàng chục là 8 ($\ge 5$) nên làm tròn lên $4800$.\\
\textbf{Đáp án của câu hỏi là C}
}
\end{ex}

\begin{ex} %[3D8H49]
Số liền trước của số $8\,000$ là
\choice
{$8001$}
{\True $7999$}
{$7990$}
{$7900$}
\loigiai{
$8000 - 1 = 7999$.\\
\textbf{Đáp án của câu hỏi là B}
}
\end{ex}

\begin{ex} %[3D8H49]
Số lớn nhất trong các số $6\,540; 6\,450; 6\,504; 6\,054$ là
\choice
{\True $6540$}
{$6450$}
{$6504$}
{$6054$}
\loigiai{
Số lớn nhất là $6540$.\\
\textbf{Đáp án của câu hỏi là A}
}
\end{ex}

\begin{ex} %[3D8V49]
Dùng các chữ số $3, 0, 8, 5$ để lập số tròn nghìn duy nhất là
\choice
{$3000$}
{$8000$}
{$5000$}
{\True Không lập được (chỉ có 1 chữ số 0, nếu lập số 4 chữ số từ các chữ số trên)}
\loigiai{
Lập số bé nhất có 4 chữ số khác nhau: $3058$.\\
\textbf{Đáp án của câu hỏi là D}
}
\end{ex}

\begin{ex} %[3D8V49]
Tính giá trị: $\text{XIX} - \text{VII} = \dots$
\choice
{\True $\text{XII}$ (số $12$)}
{$\text{XI}$}
{$\text{XIII}$}
{$\text{X}$}
\loigiai{
$19 - 7 = 12 = \text{XII}$.\\
\textbf{Đáp án của câu hỏi là A}
}
\end{ex}

\begin{ex} %[3D8V49]
Một quầy hàng có $2\,450\text{ kg}$ gạo, đã bán $450\text{ kg}$. Số gạo còn lại là
\choice
{$1900\text{ kg}$}
{\True $2000\text{ kg}$}
{$2100\text{ kg}$}
{$2050\text{ kg}$}
\loigiai{
$2450 - 450 = 2000\text{ kg}$.\\
\textbf{Đáp án của câu hỏi là B}
}
\end{ex}

\begin{ex} %[3D8V49]
Làm tròn số $6\,438$ đến hàng chục ta được
\choice
{$6430$}
{$6500$}
{\True $6440$}
{$6400$}
\loigiai{
Hàng đơn vị là 8 nên được $6440$.\\
\textbf{Đáp án của câu hỏi là C}
}
\end{ex}

\begin{ex} %[3D8V49]
Đồng hồ chỉ mấy giờ khi kim ngắn chỉ số $\text{IV}$ và kim dài chỉ số $\text{XII}$?
\choice
{\True $4$ giờ}
{$12$ giờ}
{$5$ giờ}
{$6$ giờ}
\loigiai{
Kim ngắn chỉ số 4 ($\text{IV}$), kim dài chỉ số 12 ($\text{XII}$) là 4 giờ đúng.\\
\textbf{Đáp án của câu hỏi là A}
}
\end{ex}

\begin{ex} %[3D8C49]
Tìm số có 4 chữ số, biết chữ số hàng nghìn là 5, các chữ số còn lại đều bằng nhau và tổng 4 chữ số bằng 11.
\choice
{$5111$}
{$5333$}
{$5000$}
{\True $5222$}
\loigiai{
Tổng 3 chữ số còn lại: $11 - 5 = 6$. Mỗi chữ số là $6 : 3 = 2$. Số đó là $5222$.\\
\textbf{Đáp án của câu hỏi là D}
}
\end{ex}

\textbf{Phần 2. Trắc nghiệm đúng sai}
\setcounterpageref{ex}{0}

\begin{ex} %[3D8H49]
Xét tính đúng sai về số học và chữ số La Mã:
\choiceTF
{\True Số $10\,000$ có 5 chữ số.}
{\True Chữ số La Mã $\text{XX}$ biểu thị số 20.}
{Làm tròn $650$ đến hàng trăm được $600$.}
{\True Làm tròn $650$ đến hàng trăm được $700$.}
\loigiai{
\begin{itemize}
    \item $10000$ có 5 chữ số: Đúng.
    \item $\text{XX} = 20$: Đúng.
    \item $650$ làm tròn thành 700. Ý 3 Sai, Ý 4 Đúng.
\end{itemize}
\begin{itemize}
    \item a) \textbf{Đúng}.
    \item b) \textbf{Đúng}.
    \item c) \textbf{Sai}.
    \item d) \textbf{Đúng}.
\end{itemize}
}
\end{ex}

\begin{ex} %[3D8V49]
Cho các số: $5\,670; 5\,760; 5\,607$. Xét tính đúng sai:
\choiceTF
{\True Số lớn nhất là $5\,760$.}
{\True Số bé nhất là $5\,607$.}
{Làm tròn số $5\,670$ đến hàng trăm được $5\,600$.}
{\True Làm tròn số $5\,670$ đến hàng trăm được $5\,700$.}
\loigiai{
\begin{itemize}
    \item Lớn nhất: $5760$. Đúng.
    \item Bé nhất: $5607$. Đúng.
    \item Hàng chục là 7 nên làm tròn thành 5700. Ý 3 Sai, Ý 4 Đúng.
\end{itemize}
\begin{itemize}
    \item a) \textbf{Đúng}.
    \item b) \textbf{Đúng}.
    \item c) \textbf{Sai}.
    \item d) \textbf{Đúng}.
\end{itemize}
}
\end{ex}

\textbf{Phần 3. Trả lời ngắn}
\setcounterpageref{ex}{0}

\begin{ex} %[3D8H49]
Viết số liền trước của số $6\,000$.
\shortans{5999}
\loigiai{
\textbf{Đáp số:} $5999$\\
$6000 - 1 = 5999$.
}
\end{ex}

\begin{ex} %[3D8H49]
Chữ số La Mã $\text{XIV}$ có giá trị là bao nhiêu?
\shortans{14}
\loigiai{
\textbf{Đáp số:} $14$\\
14.
}
\end{ex}

\begin{ex} %[3D8V49]
Làm tròn số $3\,648$ đến hàng trăm được số nào?
\shortans{3600}
\loigiai{
\textbf{Đáp số:} $3600$\\
$3600$.
}
\end{ex}

\begin{ex} %[3D8V49]
Tìm số lớn nhất có 4 chữ số khác nhau lập từ các chữ số $6, 0, 9, 2$.
\shortans{9620}
\loigiai{
\textbf{Đáp số:} $9620$\\
$9620$.
}
\end{ex}

\begin{ex} %[3D8V49]
Tính tổng: $\text{IX} + \text{VI}$ bằng số tự nhiên bao nhiêu?
\shortans{15}
\loigiai{
\textbf{Đáp số:} $15$\\
$9 + 6 = 15$.
}
\end{ex}

\begin{ex} %[3D8C49]
Tìm số tròn chục lớn nhất có 4 chữ số khác nhau mà chữ số hàng nghìn là 8.
\shortans{8970}
\loigiai{
\textbf{Đáp số:} $8970$\\
Để lớn nhất ta chọn hàng trăm là 9, hàng chục là 7, hàng đơn vị là 0. Số đó là $8970$.
}
\end{ex}
